%% ch01_overview.tex -- Intelligence Engineering, chapter 01 "Overview of
%% Machine Learning Systems".
%% Spine transferred from Huyen, Designing Machine Learning Systems, chapter 1.
%% The subchapters, the frame titles and the frame subtitles are the book's own
%% headings and its own one-line conditions, in the book's order.
%%
%% STATE: titles, subtitles and TEMPLATE only. No body text. Every frame carries
%% the layout it is meant to be filled into, and a % TEMPLATE comment saying why
%% that layout and what belongs in each slot. Filling them is the next pass.
%%
%% Fragment of intelligence_engineering.tex. One chapter is one lecture and
%% gets exactly ONE \lecturepage, at the top. Inside it every \subsection
%% opens on the subchapter divider, which the master emits automatically via
%% \AtBeginSubsection; do not write those divider frames by hand.

\begin{refsection}

\lecturepage{I}{Overview of a Machine Learning System}{}
\section[Overview]{Overview of a Machine Learning System}

% ======================================================================
\subchaptergloss{Don't confuse the model with the machine learning system.}
\subsection{The small part and the big picture}

% TEMPLATE: withmargin + tddef. The definition anchors the whole subchapter, so
% it takes the slide's one titled box; the five keyphrases are numbered inside
% it and each gets a frame of its own below. Margin holds the caveat that
% qualifies the definition (not a magic tool, not always optimal).
\begin{frame}{Enabling the delivery of the logic}
  \begin{withmargin}
    \begin{tdmain}
      The \textbf{model} is the small part of a machine learning system. 
      With foundation models and AI engineering, that part got smaller still. 
      The essential work of a machine learning system remains and stacks around business understanding, infrastructure, design pattern and data engineering.
    \end{tdmain}
    \begin{tdmargin}
      Before learning about solving a problem with machine learning, let's first assess when it is actually a good idea to do so.  
    \end{tdmargin}
  \end{withmargin}
\end{frame}

% TEMPLATE: withmargin. Main column takes the condition and what follows from
% it; margin takes the book's worked example (the Airbnb listing and its price).
\begin{frame}{(1) Learn}
  \begin{withmargin}
    \begin{tdmain}
        The problem you are solving, requires your system to have the capability to learn.
        Learning can look different depending on the problem and its context and boundary conditions.
        You already know about supervised learning, unsupervised learning and reinforcement learning paradigms.
    \end{tdmain}
    \begin{tdmargin}
        Machine learning is an approach to (1) learn (2) complex patterns
        from (3) existing data and use these patterns to make (4)
        predictions on (5) unseen data.
    \end{tdmargin}
  \end{withmargin}
\end{frame}

\begin{frame}{(2) Pattern}
    \begin{withmargin}
      \begin{tdmain}
          Learning requires the existence of a pattern, as opposed to noise or randomness. 
          It is not trivial nor obvious whether there is a pattern in the data, and whether it is learnable.
          Complex pattern is more difficult to embed or state explicitly, this is where machine learning has an edge over classical approaches.
      \end{tdmain}
      \begin{tdmargin}
          Machine learning is an approach to (1) learn (2) complex patterns
          from (3) existing data and use these patterns to make (4)
          predictions on (5) unseen data.
      \end{tdmargin}
    \end{withmargin}
  \end{frame}


  \begin{frame}{(3) Data}
    \begin{withmargin}
      \begin{tdmain}
          Data is available, accessible or acquirable. 
          The data is of sufficient quality and quantity to support learning, (and to validate for that matter).
          If your approach does not require data, it is because you build on embeddings or foundation models, 
          which themselves are built on data in the first place, or you need to rely on data acquisition and curation as part of your system.
      \end{tdmain}
      \begin{tdmargin}
          Machine learning is an approach to (1) learn (2) complex patterns
          from (3) existing data and use these patterns to make (4)
          predictions on (5) unseen data.
      \end{tdmargin}
    \end{withmargin}
  \end{frame}

  \begin{frame}{(4) Predictions}
    \begin{withmargin}
      \begin{tdmain}
          Machine learning is about making predictions, always. Predictions can be of different types, such as 
          classification, regression, ranking, clustering, generation and on and on.
          By reframing problems in terms of predictions, we can leverage machine learning to solve them.
          Predictions are approximations by nature and therefore come with uncertainty, 
          which needs to be accounted for.

      \end{tdmain}
      \begin{tdmargin}
          Machine learning is an approach to (1) learn (2) complex patterns
          from (3) existing data and use these patterns to make (4)
          predictions on (5) unseen data.
      \end{tdmargin}
    \end{withmargin}
  \end{frame}

  \begin{frame}{(5) Unseen Data}
    \begin{withmargin}
      \begin{tdmain}
          In a closed system or domain, the data is finite and the system can be fully specified by optimization.
          However, often the domain is open and infinite, and the system needs to be able to generalize to unseen data.
          Of course, the unseen data is not completely unknown, it is drawn from the same distribution (i.i.d.) as the 
          training data and is thus able to generalize to the unseen data inside a specific domain or distribution.
      \end{tdmain}
      \begin{tdmargin}
          Machine learning is an approach to (1) learn (2) complex patterns
          from (3) existing data and use these patterns to make (4)
          predictions on (5) unseen data.
      \end{tdmargin}
    \end{withmargin}
  \end{frame}


% TEMPLATE: two by two block grid. Four peer characteristics, none subordinate
% to another, so equal boxes rather than a bullet list. The equal height group
% keeps each row level.
\begin{frame}{Catalysts}
    A catalyst is a substance that increases the rate of reaction without 
    undergoing change itself.
  \begin{columns}[T,onlytextwidth]
    \begin{column}{0.47\textwidth}
      \begin{tdblockt}[equal height group=cat]{The task is repetitive.}
        In a sense that we have a lot of data to learn a pattern from.
      \end{tdblockt}
      \vskip0.5em
      \begin{tdblockt}[equal height group=cat]{The context changes.}
        Meaning that hard coded or other solutions will be brittle and break.
      \end{tdblockt}
    \end{column}
    \begin{column}{0.47\textwidth}
      \begin{tdblockt}[equal height group=cat]{It scales.}
        The need to scale a specific prediction capability, not only a few times.
      \end{tdblockt}
      \vskip0.5em
      \begin{tdblockt}[equal height group=cat]{Being wrong is not too bad.}
        Approximations on unseen data. Failures need to be manageable.  
      \end{tdblockt}
    \end{column}
  \end{columns}
\end{frame}

% TEMPLATE: three block row plus a closing line. The book gives exactly three
% disqualifying conditions, so one box each rather than three bullets; the line
% under them is the escape hatch, decompose and use ML on a component.
\begin{frame}{When Not to Use Machine Learning}
  \begin{columns}[T,onlytextwidth]
    \begin{column}{0.31\textwidth}
      \begin{tdblockt}[equal height group=not]{It is unethical.}
        Well, that might also be a general rule. But doing stuff at scale, without human oversight, has its implications.$^*$
      \end{tdblockt}
    \end{column}
    \begin{column}{0.31\textwidth}
      \begin{tdblockt}[equal height group=not]{There's a simple solution.}
        A simple solution, heuristic, or baseline should always be the first step regardless.
      \end{tdblockt}
    \end{column}
    \begin{column}{0.31\textwidth}
      \begin{tdblockt}[equal height group=not]{It's not economically viable.}
        The value of the prediction, needs to be greater than the cost of the system, including development, deployment and maintenance.
      \end{tdblockt}
    \end{column}
  \end{columns}
  \vskip0.8em
  $^*$ More? Enroll to my Philosophy of AI course.
  \vskip0.8em
\end{frame}







\begin{frame}{}
    \begin{statement}
      \substatement{Describe two use-cases and the rationale for whether or not to use machine learning.}  
      \alert{One} Machine Learning Use-Case and \alert{One} which is not.
      \end{statement}
\end{frame}











% The next four frames map value over technical feasibility and how that map
% moved between 2018 and 2026. Order: the 2018 baseline, the Eisenhower matrix
% in the middle, the 2026 survey numbers, the labs' own measurements. Four
% rather than the asked two to three, so there is something to cut.

% TEMPLATE: two block contrast plus a dimmed line. The 2018/2019 baseline is a
% tension between two documents, the trillion-dollar map and the months-long
% deployment reality, so each takes one levelled box and the dim line reads
% them against each other.
\begin{frame}{Mapping the value, 2018}
  \begin{columns}[T,onlytextwidth]
    \begin{column}{0.47\textwidth}
      \begin{tdblockt}[equal height group=map18]{The map: value in the trillions}
        \footnotesize
        McKinsey sized \alertbf{400+} use cases across \alertbf{19} industries and nine business functions: deep learning promised \alertbf{\$3.5 to \$5.8 trillion} in value per year, the largest shares in marketing and sales plus supply chain and manufacturing.
      \end{tdblockt}
    \end{column}
    \begin{column}{0.47\textwidth}
      \begin{tdblockt}[equal height group=map18]{The ground: months per model}
        \footnotesize
        Algorithmia surveyed \alertbf{745} practitioners in October 2019: most companies needed \alertbf{8 to 90 days} to deploy a single model, 18\% needed longer than 90 days, and deployment time grew year over year.
      \end{tdblockt}
    \end{column}
  \end{columns}
  \vskip0.8em
  {\color{tddim}The value was mapped in trillions while shipping one model took a quarter: feasibility, not value, was the bottleneck.}
  \slidesources{Chui2018,Algorithmia2020}
\end{frame}

% TEMPLATE: withmargin with the Eisenhower matrix in the main column. Redrawn
% as a quadrant matrix rather than a scatter: nobody measured the feasibility
% of customer service, so plotted dots claimed a precision the data does not
% have, and eleven of them collided with each other and with the quadrant
% verbs. Now each quadrant is its verb plus the domains that sit in it. The
% frontier between hard and works is the one red element and it carries the
% claim: it moves right. Margin holds the numbers that moved it, one per
% source, so the diagram states the argument and the margin evidences it.
\begin{frame}{Value over technical feasibility}
  \begin{withmargin}
    \begin{tdmain}
      \centering
      \begin{tikzpicture}[
          ax/.style={tddim, semithick, -{Latex[length=4pt,width=3pt]}},
          split/.style={tddim, thin, densely dashed},
          quad/.style={text=tddim, font=\ttfamily\scriptsize, anchor=north west},
          dom/.style={text=tdfg, font=\tiny, anchor=north west, align=left,
                      inner sep=0pt}]
        % axes: left and bottom only, no enclosing box
        \draw[ax] (0,0) -- (0,4.40);
        \draw[ax] (0,0) -- (7.75,0);
        % the value split is a plain guide
        \draw[split] (0,2.10) -- (7.6,2.10);
        % the feasibility split is the frontier, and it moves
        \draw[DHBWred, semithick, densely dashed] (3.8,0) -- (3.8,4.25);
        \draw[DHBWred, semithick, -{Latex[length=4pt,width=3pt]}]
          (3.8,4.45) -- (4.6,4.45);
        \node[text=DHBWred, font=\tiny, anchor=west] at (4.68,4.45)
          {the frontier moves};
        % x axis: technical feasibility
        \node[tddim, font=\tiny, anchor=north west] at (0,-0.08) {hard today};
        \node[tddim, font=\tiny, anchor=north east] at (7.6,-0.08) {works today};
        \node[tddim, font=\scriptsize, anchor=north] at (3.8,-0.45)
          {technical feasibility};
        % y axis: business value
        \node[tddim, font=\tiny, rotate=90, anchor=south west] at (-0.08,0) {modest};
        \node[tddim, font=\tiny, rotate=90, anchor=south east] at (-0.08,4.25)
          {transformative};
        \node[tddim, font=\scriptsize, rotate=90, anchor=south] at (-0.52,2.10)
          {business value};
        % invest: high value, still hard
        \node[quad] at (0.20,4.20) {invest};
        \node[dom] at (0.20,3.80) {autonomous driving\\medical diagnosis\\agentic workflows end to end};
        % deploy now: high value, works today
        \node[quad] at (4.05,4.20) {deploy now};
        \node[dom] at (4.05,3.80) {coding agents\\recommender systems\\fraud detection\\machine translation};
        % skip: hard and modest
        \node[quad] at (0.20,1.95) {skip};
        \node[dom] at (0.20,1.55) {bespoke one off models};
        % automate cheaply: works today, modest value
        \node[quad] at (4.05,1.95) {automate cheaply};
        \node[dom] at (4.05,1.55) {speech transcription\\document OCR};
      \end{tikzpicture}
    \end{tdmain}
    \begin{tdmargin}
      SWE-bench Verified rose from 60\% to near 100\% within a year.
      \vskip0.5em
      Claude Opus 4.1 wins or ties against human experts on 47.6\% of GDPval tasks.
      \vskip0.5em
      Among organizations above \$1 billion revenue, 40\% now scale AI agents, up from 27\% in one year.
    \end{tdmargin}
  \end{withmargin}
  \slidesources{AIIndex2026,OpenAI2025,McKinsey2026}
\end{frame}

% TEMPLATE: side image hero plus two label led columns. Mark's call, and the
% reason holds: a pair of cards frames the two shares as peers of equal weight
% when the point is that they disagree, so the boxes come off and the column
% split alone carries the contrast. The strip takes the image; the headings
% survive as bold lead lines rather than block titles.
\begin{frame}{The field is moving, 2026}
  \sideimagecover[0.33]{placeholder_gray.png}
  % Asset brief: the adoption curve against the value capture curve on one set
  % of axes, 2023 to 2026, the two diverging. Caption stays under ~50
  % characters: it is a rotated node with no text width and the page clips
  % whatever runs past the top edge.
  \sidecaption{Adoption climbs, value capture flat}
  \begin{sidebody}
    \scriptsize
    \begin{columns}[T,onlytextwidth]
      \begin{column}{0.47\linewidth}
        \textbf{Adoption is near universal}
        \begin{itemize}
          \item \alertbf{88\%} of organizations use AI in at least one business function, \alertbf{72\%} use gen AI.
          \item Gen AI reached 53\% population adoption within three years, faster than the PC or the internet.
          \item US private AI investment reached \$285.9 billion in 2025.
        \end{itemize}
      \end{column}
      \hfill
      \begin{column}{0.47\linewidth}
        \textbf{Value capture is not}
        \begin{itemize}
          \item \alertbf{80\%} of AI users report individual productivity gains.
          \item \alertbf{37\%} of organizations attribute any EBIT impact to AI, flat since 2025.
          \item \alertbf{6\%} qualify as AI high performers.
        \end{itemize}
      \end{column}
    \end{columns}
  \end{sidebody}
  \slidesources{McKinsey2026,AIIndex2026}
\end{frame}

% TEMPLATE: two block contrast plus a dimmed line. The two labs publish the
% same map from their own vantage points, a benchmark against experts and a
% telescope on live traffic, so one box per lab and the dim line reads the
% two together.
\begin{frame}{What the labs measure}
  \begin{columns}[T,onlytextwidth]
    \begin{column}{0.47\textwidth}
      \begin{tdblockt}[equal height group=labs]{OpenAI: GDPval}
        \footnotesize
        \alertbf{1{,}320} tasks from \alertbf{44} occupations across nine sectors, graded head to head by experts averaging 14 years in the field. Frontier models finish them about \alertbf{100$\times$} faster and \alertbf{100$\times$} cheaper than the experts.
      \end{tdblockt}
    \end{column}
    \begin{column}{0.47\textwidth}
      \begin{tdblockt}[equal height group=labs]{Anthropic: Economic Index}
        \footnotesize
        \alertbf{49\%} of jobs show Claude on at least a quarter of their tasks. Augmentation leads automation \alertbf{52\% to 45\%}. The single most common task: modifying software to correct errors, 6\% of all usage.
      \end{tdblockt}
    \end{column}
  \end{columns}
  \vskip0.8em
  {\color{tddim}Broad across occupations, thin within them, and the densest use is code: the labs' own traffic redraws the feasibility axis every quarter.}
  \slidesources{OpenAI2025,Anthropic2026}
\end{frame}











% ======================================================================
\subchaptergloss{Research and Production are not the same.}
\subsection{Understanding Machine Learning Systems}

% TEMPLATE: comparison table. Table 1-1 of the book is already a table and
% should stay one: five rows, research against production. booktabs only, no
% vertical rules. The five row labels are the book's.
\begin{frame}{Horizontal versus Vertical}
  \footnotesize
  Nothing motivates \alert{system} thinking more than the need to go \alert{horizontal}.
  Academia is vertical, with a narrow focus on a single problem, where drilling deep, fast is rewarded and
  focus allows to teach a wide range of concepts.
  \vspace{1.5em}

  \renewcommand{\arraystretch}{1.25}
  
  \begin{tabular}{@{}%
    >{\raggedright\arraybackslash}p{0.22\textwidth}%
    >{\raggedright\arraybackslash}p{0.35\textwidth}%
    >{\raggedright\arraybackslash}p{0.35\textwidth}@{}}
    \toprule
     & \textbf{Research} & \textbf{Production}\\
    \midrule
    \textbf{Objectives} & State of the Art highly engineered models &  Stakeholder outcome driven \\
    \textbf{Compute} & Training and high throughput &  Fast inference with low latency \\
    \textbf{Data} & Often static benchmark oriented & Shifting domains and structures  \\
    \textbf{Usability metrics} & Usually not considered & UI, Fairness, Robust, Interpretability \\
    \bottomrule
  \end{tabular}
\end{frame}









% TEMPLATE: side image hero plus a two row card grid. The five stakeholders are
% peers, none subordinate to another, so a grid rather than a stack, and two
% rows keep the sidebody inside its height budget. Each card carries an
% objective and the solution space its owner already has in mind, which are the
% two axes of the Eisenhower map in the strip. No red on the cards: the one red
% element on this frame is the term the whole slide is about.
\begin{frame}{Stakeholders and Objectives}
  \sideimagecover[0.33]{placeholder_gray.png}
  % Asset brief: an Eisenhower map, concern or objective over solution space,
  % both axes running from business to technical. The visible caption has to
  % stay under ~50 characters: it is a rotated node with no text width and
  % the page clips whatever runs past the top edge.
  \sidecaption{Eisenhower: objective over solution space}
  \begin{sidebody}
    \scriptsize
    Tinder for horses. In research the objective is one number on one benchmark.
    Ship the same model as a product and five people want five different things
    from it, each arriving with a solution space already in mind:
    \alertbf{conflicting requirements} are the norm, and not all of them are must haves.
    \vskip0.7em
    \tiny
    \begin{columns}[T,onlytextwidth]
      \begin{column}{0.315\linewidth}
        \begin{tdblockt}[coltitle=tdfg,before upper=\raggedright,equal height group=sh,left=4pt,right=4pt,toptitle=1pt,top=1pt,bottom=2pt]{ML engineers}
          {\color{tddim}Objective} better matches.\par
          {\color{tddim}Solution space} more data, a bigger model.
        \end{tdblockt}
      \end{column}
      \hfill
      \begin{column}{0.315\linewidth}
        \begin{tdblockt}[coltitle=tdfg,before upper=\raggedright,equal height group=sh,left=4pt,right=4pt,toptitle=1pt,top=1pt,bottom=2pt]{Platform team}
          {\color{tddim}Objective} nothing pages at 3am.\par
          {\color{tddim}Solution space} freeze the model, build the platform.
        \end{tdblockt}
      \end{column}
      \hfill
      \begin{column}{0.315\linewidth}
        \begin{tdblockt}[coltitle=tdfg,before upper=\raggedright,equal height group=sh,left=4pt,right=4pt,toptitle=1pt,top=1pt,bottom=2pt]{Product team}
          {\color{tddim}Objective} the swipe feels instant.\par
          {\color{tddim}Solution space} a latency budget of 100 ms.
        \end{tdblockt}
      \end{column}
    \end{columns}
    \vskip0.35em
    \begin{columns}[T,onlytextwidth]
      \begin{column}{0.315\linewidth}
        \begin{tdblockt}[coltitle=tdfg,before upper=\raggedright,equal height group=sh2,left=4pt,right=4pt,toptitle=1pt,top=1pt,bottom=2pt]{Sales team}
          {\color{tddim}Objective} revenue per match.\par
          {\color{tddim}Solution space} rank the premium stallions first.
        \end{tdblockt}
      \end{column}
      \hfill
      \begin{column}{0.315\linewidth}
        \begin{tdblockt}[coltitle=tdfg,before upper=\raggedright,equal height group=sh2,left=4pt,right=4pt,toptitle=1pt,top=1pt,bottom=2pt]{Manager}
          {\color{tddim}Objective} margin.\par
          {\color{tddim}Solution space} fewer models, and maybe fewer engineers.
        \end{tdblockt}
      \end{column}
      \hfill
      \begin{column}{0.315\linewidth}
      \end{column}
    \end{columns}
  \end{sidebody}
  \slidesources{Huyen2022}
\end{frame}


% TEMPLATE: side image hero plus a bullet list. Three separate charges against
% the same practice, so a bullet each rather than boxes; none of them is a
% comparison, so no paired blocks. The strip takes the joke the criticism is
% named after. Picard's seed scan is the whole argument of the second bullet,
% so its numbers sit in the bullet and not in a caption.
\begin{frame}{Leaderboards and Benchmarks}
  \sideimagecover[0.33]{placeholder_gray.png}
  % Asset brief: a trust me bro benchmark chart, bars with no axis, no
  % baseline and no error bars, posted to announce a model. Caption kept
  % under ~50 characters, see the note on the frame above.
  \sidecaption{Trust me bro benchmark: no axis, no error bars}
  \begin{sidebody}
    \scriptsize
    A leaderboard scores the \alertbf{model}. A system is not a model.
    \vskip0.5em
    \begin{itemize}
      \item What a benchmark measures is one metric on one frozen test set. What it
        never measures is the pipeline around it: acquisition, cleaning, labelling,
        feature engineering, deployment, monitoring.
      \item Thousands of entries against the same test set is multiple hypothesis
        testing, so part of the podium is pure chance. Scanning 10,000 seeds on
        CIFAR 10 spreads accuracy from 89.01 to 90.83 percent, 1.82 points, with
        nothing changed but the seed.
      \item A benchmark is a proxy for a single metric. It was never forged in
        production, where objectives conflict, data shifts, and the users are real.
    \end{itemize}
  \end{sidebody}
  \slidesources{Picard2021}
\end{frame}

\begin{frame}{Computational priorities}
  \begin{columns}[T,onlytextwidth]
    \begin{column}{0.47\textwidth}
      \begin{tdblockt}[equal height group=cp]{Research: fast training}
        \scriptsize
        \begin{itemize}
          \item Training is the bottleneck: many models, each doing multiple passes over the data.
          \item Prioritizes \alertbf{high throughput}, how many queries are processed in a period of time.
        \end{itemize}
      \end{tdblockt}
    \end{column}
    \begin{column}{0.47\textwidth}
      \begin{tdblockt}[equal height group=cp]{Production: fast inference}
        \scriptsize
        \begin{itemize}
          \item Once deployed, inference is the bottleneck: the model's job is to generate predictions.
          \item Prioritizes \alertbf{low latency}, the time from receiving a query to returning the result.
        \end{itemize}
      \end{tdblockt}
    \end{column}
  \end{columns}
  \vskip0.3em
  \scriptsize
  {\centering
    \begin{tabular}{@{}lll@{}}
      \toprule
      \textbf{Processing} & \textbf{Average latency} & \textbf{Throughput}\\
      \midrule
      One query at a time & 10 ms & 100 queries/second\\
      One query at a time & 100 ms & 10 queries/second\\
      Batches of 10 queries & 10 ms & 1,000 queries/second\\
      Batches of 50 queries & 20 ms & 2,500 queries/second\\
      \bottomrule
    \end{tabular}\par}
  \vskip0.3em
  {\color{tddim}Batched, higher latency might also mean higher throughput.}
  \slidesources{Huyen2022}
\end{frame}

% TEMPLATE: withmargin plus a takeaway. The main column runs the book's ten
% measured latencies and its percentile argument in order, because the point is
% arithmetic. Margin holds the three dated production studies as the evidence
% that latency matters. The takeaway is the Amazon observation, the reason the
% tail of the distribution is not a rounding error.
\begin{frame}{Latency is a distribution}
  \begin{withmargin}
    \begin{tdmain}
      \small
      Ten requests, latency in ms: 100, 102, 100, 100, 99, 104, 110, 90, \alertbf{3,000}, 95.

      The average latency is 390 ms, which makes your system seem slower than it actually is.

      Better to think in percentiles: p50 is the median, p90 and p99 expose the outliers.

      The p90 here is 3,000 ms, the one request to investigate.
    \end{tdmain}
    \begin{tdmargin}
      Latency matters a lot.
      \vskip0.5em
      2017, Akamai: a 100 ms delay can hurt conversion rates by 7\%.
      \vskip0.5em
      2019, Booking.com: about 30\% more latency cost about 0.5\% in conversion rates, ``a relevant cost for our business''.
      \vskip0.5em
      2016, Google: more than half of mobile users leave a page that takes more than three seconds to load.
    \end{tdmargin}
  \end{withmargin}
  \begin{takeaway}
    On Amazon the slowest requests come from the most valuable customers.
  \end{takeaway}
  \slidesources{Huyen2022}
\end{frame}

% TEMPLATE: two block contrast. The book's Data subheading is a straight
% research against production split, so each world takes one box, headed with
% its own entry from Table 1-1 and filled with its data properties in the
% book's words. The dimmed line under them is the book's Karpathy anchor for
% the same contrast.
\begin{frame}{Data}
  \begin{columns}[T,onlytextwidth]
    \begin{column}{0.47\textwidth}
      \begin{tdblockt}[equal height group=data]{Research: static}
        \footnotesize
        \begin{itemize}
          \item Clean and well formatted
          \item Static, so the community can benchmark
          \item Quirks of the dataset are known
          \item Open source scripts to feed your models
          \item Historical data, already stored somewhere
        \end{itemize}
      \end{tdblockt}
    \end{column}
    \begin{column}{0.47\textwidth}
      \begin{tdblockt}[equal height group=data]{Production: constantly shifting}
        \footnotesize
        \begin{itemize}
          \item Messy, noisy, possibly unstructured
          \item \alert{Likely biased, and you do not know how}
          \item Labels sparse, imbalanced, or incorrect
          \item Changing requirements can mean relabeling
          \item Privacy and regulatory concerns
          \item Generated by users, systems, third parties
        \end{itemize}
      \end{tdblockt}
    \end{column}
  \end{columns}
  \vskip0.5em
  {\color{tddim}\footnotesize Andrej Karpathy, director of AI at Tesla, drew the data problems of his PhD against those of his time at Tesla.}
  \slidesources{Huyen2022}
\end{frame}

% TEMPLATE: one quote block between two lines. The book's fairness point is a
% quotation and its rebuttal rather than a comparison, so the deferral gets the
% slide's one box and the verdict lands inside it. The line above is the
% research phase condition, the line below the asymmetry, and the takeaway the
% leaderboard that does not exist.
\begin{frame}{Fairness}
  During the research phase, a model is not yet used on people, so it is easy for researchers to put off fairness as an afterthought.
  \vskip0.4em
  \begin{tdblock}
    {\color{tddim}``Let's try to get state of the art first and worry about fairness when we get to production.''}
    \vskip0.5em
    \alertbf{When it gets to production, it is too late.}
  \end{tdblock}
  \vskip0.4em
  Optimize for better accuracy or lower latency and you can show that you beat state of the art.
  \begin{takeaway}
    As of writing the book, there is no equivalent state of the art for fairness metrics.
  \end{takeaway}
  \slidesources{Huyen2022}
\end{frame}

% TEMPLATE: withmargin with a bullet list. The book argues this one by naming
% victims, so the main column is one book example per bullet, opened by the
% line that explains them all and closed by what deployment at scale does to
% them. Margin takes the sourced evidence, the Berkeley lending finding and the
% McKinsey number that says how few companies act on it.
\begin{frame}{Fairness at scale}
  \begin{withmargin}
    \begin{tdmain}
      \small
      Machine learning algorithms do not predict the future, they encode the past, thus perpetuating the biases in the data.
      \begin{itemize}
        \item Your loan application might be rejected because the algorithm picks on your \alert{zip code}, which embodies biases about one's socioeconomic background.
        \item Your resume might be ranked lower because the ranking system employers use picks on the \alert{spelling of your name}.
        \item Your mortgage might get a higher interest rate because it relies partially on \alert{credit scores}, which favor the rich and punish the poor.
      \end{itemize}
      \alertbf{Deployed at scale, they discriminate at scale}, judging millions in split seconds.
    \end{tdmain}
    \begin{tdmargin}
      In 2019, Berkeley researchers found that both face-to-face and online lenders rejected \textbf{1.3 million} creditworthy Black and Latino applicants between 2008 and 2015. With the income and credit scores kept and the race identifiers deleted, the mortgage application was accepted.
      \vskip0.5em
      McKinsey \& Company, 2019: only \textbf{13\%} of the large companies surveyed said they are taking steps to mitigate risks to equity and fairness.
    \end{tdmargin}
  \end{withmargin}
  \slidesources{Huyen2022}
\end{frame}

% TEMPLATE: two block contrast plus a dimmed closing line. Hinton's question is
% a forced choice between two options, so each option takes one box and the two
% cure rates read off against one another. The lines above set the scenario, the
% quoted question sits under the boxes, and the dimmed line is what a room of
% executives answered when the book's author put the question to them.
\begin{frame}{Interpretability}
  In early 2020, Turing Award winner Geoffrey Hinton posed a heatedly debated question.
  Suppose you have cancer and you have to choose.
  \vskip0.9em
  \begin{columns}[T,onlytextwidth]
    \begin{column}{0.47\textwidth}
      \begin{tdblockt}[equal height group=surgeon]{Black box AI surgeon}
        \alertbf{90\%} cure rate, and it cannot explain how it works.
      \end{tdblockt}
    \end{column}
    \begin{column}{0.47\textwidth}
      \begin{tdblockt}[equal height group=surgeon]{Human surgeon}
        \alertbf{80\%} cure rate.
      \end{tdblockt}
    \end{column}
  \end{columns}
  \vskip0.9em
  \alertbf{``Do you want the AI surgeon to be illegal?''}
  \vskip0.6em
  {\color{tddim}Put to a group of 30 technology executives at public nontech companies, half of them wanted the AI surgeon to operate on them. The other half wanted the human surgeon.}
  \slidesources{Huyen2022}
\end{frame}

% TEMPLATE: withmargin. The main column carries the asymmetry this subchapter
% is built on, research optimizes one objective and industry owes an
% explanation, with the book's two reasons as bullets. Margin holds why the
% discomfort is specific to AI, and the number showing that a requirement is
% not yet a practice.
\begin{frame}{Interpretability in production}
  \begin{withmargin}
    \begin{tdmain}
      \small
      Most machine learning research is still evaluated on a single objective, model performance, so researchers are not incentivized to work on model interpretability.
      In industry, interpretability is not just optional for most use cases but a requirement.
      \begin{itemize}
        \item Important for \alertbf{users}, both business leaders and end users, to understand why a decision is made, so that they can trust a model and detect potential biases.
        \item Important for \alertbf{developers}, to be able to debug and improve a model.
      \end{itemize}
    \end{tdmain}
    \begin{tdmargin}
      Most of us are comfortable using a microwave without understanding how it works. Many do not feel the same way about AI, especially if that AI makes important decisions about their lives.
      \vskip0.5em
      As of 2019, only \textbf{19\%} of large companies are working to improve the explainability of their algorithms.
    \end{tdmargin}
  \end{withmargin}
  \slidesources{Huyen2022}
\end{frame}


% TEMPLATE: single block plus a takeaway. The book closes the research versus
% production subchapter with an argument and its rebuttal rather than a
% comparison, so the claim takes the slide's one box and the takeaway carries
% what the book answers. Neither a margin nor a pair of boxes, because both
% neighbours already use those.
\begin{frame}{Discussion}
  \begin{tdblock}
  \end{tdblock}
  \begin{takeaway}
  \end{takeaway}
  \slidesources{Huyen2022}
\end{frame}

% TEMPLATE: two block contrast plus a takeaway. Departs from the stub's
% withmargin on purpose: the frame before it is a withmargin and a second in a
% row stalls, and the book's point here is symmetric, what software engineering
% assumes against what a machine learning system is actually made of, which
% reads as two levelled boxes rather than a claim with an aside. Each box holds
% one world in the book's own properties; the takeaway is the book's own answer
% to the question the subheading opens with.
\begin{frame}{Machine Learning Systems Versus Traditional Software}
  \begin{columns}[T,onlytextwidth]
    \begin{column}{0.47\textwidth}
      \begin{tdblockt}[equal height group=mlsw]{Traditional software engineering}
        \scriptsize
        \begin{itemize}
          \item An underlying assumption that \alertbf{code and data are separated}.
          \item Keep things as modular and separate as possible: separation of concerns.
          \item Many of its tools do work for machine learning.
        \end{itemize}
      \end{tdblockt}
    \end{column}
    \begin{column}{0.47\textwidth}
      \begin{tdblockt}[equal height group=mlsw]{Machine learning systems}
        \scriptsize
        \begin{itemize}
          \item \alertbf{Part code, part data}, and part artifacts created from the two.
          \item Applications with the most and best data win: companies improve data, not algorithms.
          \item Data can change quickly, so applications need to be adaptive.
        \end{itemize}
      \end{tdblockt}
    \end{column}
  \end{columns}
  \begin{takeaway}
    ML production would be a much better place if ML experts were better software engineers.
  \end{takeaway}
  \slidesources{Huyen2022}
\end{frame}

% TEMPLATE: withmargin, the layout this chapter uses when one claim needs one
% worked example to land, and the stub's intent honoured one frame later. Main
% column runs the book's argument in its order: code only, then code and data,
% then the two open questions, then the inequality of samples, closing on the
% risk that follows. Margin holds the book's lung scan counts, which are the
% whole proof that samples are not equal.
\begin{frame}{Test and version the data too}
  \begin{withmargin}
    \begin{tdmain}
      \small
      In traditional software engineering, you only need to focus on testing and versioning your code.
      With machine learning, we have to test and version our data too, and \alertbf{that is the hard part}.
      \begin{itemize}
        \item How to version large datasets?
        \item How to know if a data sample is good or bad for your system?
      \end{itemize}
      Not all data samples are equal. Some are more valuable to your model than others.
      \vskip0.3em
      \alertbf{Indiscriminately accepting all available data might hurt performance and make your model susceptible to data poisoning attacks.}
    \end{tdmain}
    \begin{tdmargin}
      Your model has already trained on \textbf{one million} scans of normal lungs and only \textbf{one thousand} scans of cancerous lungs.
      \vskip0.5em
      A scan of a cancerous lung is much more valuable than a scan of a normal lung.
    \end{tdmargin}
  \end{withmargin}
  \slidesources{Huyen2022}
\end{frame}

% TEMPLATE: three block row, two dimmed asides, one closing band. The book
% names three engineering challenges of the same kind, size, speed and
% visibility, so one box each rather than a bullet list, levelled by the equal
% height group. Under the row sit the book's own two one-liners, the joke about
% how fast this ages and the autocompletion example that says what fast enough
% means. The band breaks the row on purpose: it is the last thing the chapter
% says, so the good news gets its own box and keeps the BERT numbers that were
% once called impractical.
\begin{frame}{Size, speed and visibility}
  \footnotesize
  \begin{columns}[T,onlytextwidth]
    \begin{column}{0.31\textwidth}
      \begin{tdblockt}[equal height group=ssv]{Size}
        \scriptsize
        As of 2022 it is common for models to have \alertbf{hundreds of millions, if not billions, of parameters}, needing gigabytes of RAM.
      \end{tdblockt}
    \end{column}
    \begin{column}{0.31\textwidth}
      \begin{tdblockt}[equal height group=ssv]{Speed}
        \scriptsize
        Getting large models into production, \alertbf{especially on edge devices}, is a massive challenge, as is running them fast enough.
      \end{tdblockt}
    \end{column}
    \begin{column}{0.31\textwidth}
      \begin{tdblockt}[equal height group=ssv]{Visibility}
        \scriptsize
        Monitoring and debugging in production is also nontrivial: models get complex, with a \alertbf{lack of visibility} into their work.
      \end{tdblockt}
    \end{column}
  \end{columns}
  \vskip0.3em
  {\color{tddim}\scriptsize
    An autocompletion model is useless if the time it takes to suggest the next character is longer than the time it takes for you to type.}
  \vskip0.3em
  \begin{tdblockt}{The good news: these challenges are being tackled at a breakneck pace}
    \scriptsize
    In 2018 people said BERT was too big, too complex and too slow to be practical. The pretrained large BERT model has \alertbf{340 million} parameters and is \alertbf{1.35 GB}. Two years later, BERT and its variants were used in almost every English search on Google.
  \end{tdblockt}
  \slidesources{Huyen2022}
\end{frame}


% ----------------------------------------------------------------------
% This chapter's own references. The refsection opened above scopes the
% citation numbering to this chapter, so chapter_pdfs/<this>.pdf resolves
% its own [n] markers instead of pointing at a list it does not contain.
\begin{frame}{References}
  \printbibliography[heading=none]
\end{frame}
\end{refsection}
